\cotefile{15}
\scanpage{34}

% B1 | text | head and first five lines
\uline{Ex. 3} Soit fixé, \unsure{un|une} \unsure{fibre|filtre} de la catégorie des
\unsure{Dessin|Dessins}$_{\otimes}$ $C$ sur $A$ comme dans Ex. 2, avec
élément $\tau$ de $V_A(C)$. Si $E$ est \inserted{\unsure{centrale|centré}} de rang $n$,
il faut considérer que les $C$ \uline{généralisent} \inserted{les \deleted{(\unsure{de même si on se restreint})}}
\unsure{des|les} liens $\unsure{G|G_m}$ \uuline{sur $A$} trivialisés par \unsure{$K$|un},

% B2 | text | fine writing between lines five and six
(qui correspond au cas où \deleted{\unsure{sur}}, dans \unsure{tt}, ordre maximal est \deleted{\unsure{un}}
\deleted{et \unsure{donc} (cas où $E$ est \unsure{déployée} sur \unsure{tt} complété de \unsure{$A$})}
une algèbre d'Azumaya sur $R$, i.e. la classe de \unsure{lien} $\mathrm{Br}(K)$ est dans
l'image de $\mathrm{Br}(R)$)

% B3 | text | left margin opposite opening
% slanted
\margin{gerbe sur $A$ de lien $G_m$ \inserted{et trivialisée par \unsure{$K$}} $\longmapsto$ catégorie $C$ des représentations de cette gerbe qui sont de degré $n$, poids $1$}

% B4 | diagram | left margin below the first marginal note
% slanted
\margin{\diagram{\(\begin{array}{c@{\quad}l}
 C \longmapsto & \text{gerbe des}\\[-2pt]
 \text{\scriptsize cat. de} & \text{foncteurs}\\
 \text{\scriptsize \unsure{Dessin}}_{\otimes} & \text{admissibles}\\
 \text{\scriptsize \unsure{munie} $\tau$} & \text{de $C$ dans les}\\
 \text{\scriptsize où $\{\tau\}\in$} & \text{modules \unsure{sur}}\\
 \operatorname{Im}(\mathrm{Br}\,A\to\mathrm{Br}\,K) & \text{(trivialisée \unsure{par}}\\
 & \text{\unsure{de} $C$)}
\end{array}\)}}

% B5 | text | continuation through middle of sheet
et il \deleted{faut} voir de \unsure{$=$|même} que les couples $(C,\tau)$
généralisent la donnée d'une gerbe de groupe $\mu_n$.

% B6 | text | middle, new paragraph
Pour ceci, \deleted{\unsure{on}} notons \deleted{qu'\unsure{une} \unsure{gerbe} \inserted{la catégorie des}} qu'une
gerbe de lien $\mu_n$ sur $A$ définit bien une $(C,\tau)$ :
$C$ : les \deleted{\unsure{trivialisations} de} représentations de degré $n$, de rang $1$
de la gerbe, $\tau=$ \deleted{la représentation déterminant}

% B7 | text | middle, new paragraph beginning Nous
Nous allons déterminer d'autre part les autoéquivalences de $(\unsure{E|C},\tau)$ \unsure{;}
en termes d'un objet $X$ de $C$, \inserted{d'où $O=\operatorname{End}(X)$, $\nu$}
c'est la catégorie des \deleted{\unsure{objets}} bimodules
$P$ sous $O$, inv. loc. sur $A$ des deux cotés,
munis d'un isom \inserted{dans $\mathbb P(A)$} ($\deleted{\nu(X)=L\otimes}\;\tau=L\otimes_A\nu(X)$)

% B8 | display | lower middle, centred norm
\[
 N_{O/A}(P_O)\simeq\deleted{\unsure{O}} A.
\]

% B9 | text | left margin of norm paragraph, braced
\margin{$\varphi(X)=P\otimes_O X$\newline
$\nu(\varphi(X))=N(P)\otimes\nu(X)$\newline
$\phantom{\nu(\varphi(X))}\unsure{\wr\!|\,|\,\simeq}$\newline
$\varphi(\nu(X))$}

% B10 | text | below norm and outlined cancellation
\deleted{compatible avec}.
\deleted{Supposons que tous les \unsure{$P$} \inserted{de \unsure{$\mathbb P(O/A)$}} soient isomorphes
ce qui est le cas si $A$ est local, ou semi local. Alors}

% B11 | text | foot of sheet
C'est une \unsure{Gr.|Gp.}-catégorie \inserted{et $=$ une catégorie de Picard strt. \uline{comm.}}, où les endom.
de l'objet $1$ sont les éléments $\lambda\in A^*$ tels que $\lambda^n=1$ ;
\deleted{\unsure{cette}} catégorie de Picard contient comme sous-catégorie pleine $\underline{\mathrm{Tors}}(S,\mu_n)$ des
% archive sheet number: 33

\scanpage{36}

% B1 | text | opening line, continuing scan 34
torseurs \deleted{\unsure{sur}} sur $S$ de groupe $\mu_n$, et son $\pi_1$ \deleted{\unsure{est}} s'insère ds

% B2 | diagram | top horizontal sequence
\diagram{\begin{tikzcd}[ampersand replacement=\&,column sep=small]
0\arrow[r]\& \unsure{H^1|H}(S,\mu_n)\arrow[r]\& \pi_0(\underline{\unsure{\mathrm{Bordt}|\mathrm{Bord}\tau}}(C,\tau))\arrow[r]\& \deleted{\unsure{\mathrm{Br}| \mathrm{Pic}}}\inserted{\mathrm{Ram}(C/A)}\arrow[r]\&0
\end{tikzcd}}

% B3 | text | below sequence
\deleted{(Si $\mathrm{Pic}(A)=0$, \unsure{cette catégorie}}

% B4 | display | definition and note on the summand
\[
\text{où}\qquad \mathrm{Ram}(C/A)\overset{\mathrm{déf}}{=}\deleted{\unsure{\sum_q}}\mathrm{Bordt}(C)/\mathrm{Pic}(A)
\simeq\sum_p\mathbb Z/e_p(C)\,\mathbb Z.
\]
\margin{degré de ramification}

% B5 | text | upper middle, new paragraph
\deleted{Mais} \deleted{Dans le cas} Soit maintenant $A'$ une $A$-algèbre \unsure{(comm.)}
et $C'$ une catégorie de \unsure{Dessin}$_\otimes$ sur $A'$, \inserted{et $\tau'\in\unsure{V}(C')$,}
supposons $(C_{A'},T_{A'})\underset{A'}{\simeq}(C',\unsure{A'|T'})$
\deleted{Alors l'ensemble des}
(il suffit pour ceci que $C_{A'}\simeq C'$ et $\mathrm{Pic}(A')=0$,
p. ex. $A'$ local) Alors \unsure{l'ens.|l'une} des équivalences
de $(C_{A'},T_{A'})$ avec $(C',T')$, i.e. des couples

% B6 | display | middle, specification in two lines
\[
\begin{aligned}
\text{$(\varphi,u)$ avec $\varphi$}\quad &\deleted{\text{équivalence}}\quad C_{A'}\mathrel{\overset{\varphi}{\unsure{-|\longrightarrow}}} C'\quad \deleted{\unsure{\mathrm{et}}}\quad A\text{-linéaire adm.}\\
&\hspace{5em}\text{induisant }\simeq C_{A'}\to C',\text{ et}\\
u:\quad &\nu(\varphi)(T_{A'})\overset{\sim}{\longrightarrow}T'\quad ;
\end{aligned}
\]

% B7 | text | lower middle, the two actions
est une \uline{catégorie} \uline{torseur} sous la \inserted{$\unsure{Gr.|Gr}$} catégorie
\unsure{$\mathrm{gr}$} $\underline{\mathrm{Equ}}(C_{A'},T_{A'})$ (à droite) ou $\underline{\mathrm{Equ}}(C',T')$ (à gauche)

% B8 | text | bottom paragraph and marginal check
Si p. ex. $\mathrm{Ram}(C'/A')=0$ ($C'$ non ramifiée sur $A'$)
\deleted{\unsure{on aura}} alors c'est une catégorie-torseur sous
\deleted{\unsure{$H^1$}} $\underline{\underline{\mathrm{Tors}}}(S',\mu_n)$. En fait, \unsure{on a} une équivalence
\uline{canonique} $\underline{\mathrm{Equ}}(C_{A'},T_{A'})\simeq\underline{\mathrm{Equ}}(C',T')$ (car on a
une \deleted{\unsure{gr.}} catégorie \uuline{de Picard}) et par suite
% slanted
\margin{\uline{À vérifier}!!}
les deux structures de \unsure{tors.}-catégories peuvent
s'identifier\ldots
% archive sheet number: 35

\scanpage{38}

% B1 | text | first five lines
Supposons qu'on ait pu associer \uline{canoniqt},
à la catégorie de \unsure{Dessin|Dessins}$_\otimes$-Brauer $C$ (en termes de
la seule classe d'équivalence de $C$), un objet
$T$ de $\unsure{\mathcal V|V}(C)$. Cela signifie en particulier que
nous avons relevé le foncteur canonique \inserted{de Gr-cat}

% B2 | diagram | upper pair of opposite arrows
\diagram{\begin{tikzcd}[ampersand replacement=\&]
\underline{\mathrm{Equ}}(C,T)\arrow[r]\&\underline{\mathrm{Equ}}(C)\arrow[l,bend left=27,"{\deleted{\unsure{\pi}}\,c}"]
\end{tikzcd}}
\unsure{en}

% B3 | text | below curved arrow
Donc pour tte catégorie $\unsure{C_1|C'}$ équivalente à $C$,
on obtient un objet $T_{\unsure{1|C'}}^{\unsure{c}}$ canonique dans $\nu(C')$.

% B4 | text | middle through quotient
On voit donc que pour tt $(A',C',T')$ comme
plus haut, \deleted{l'ens. des} tel que $(C',T')$ soit
équivalent à $(C,T)_{A'}$, les équivalences
en question font une catégorie-torseur sous
$\underline{\mathrm{Equ}}((C,T)_{A'})$ \inserted{$\simeq\underline{\mathrm{Equ}}(C',T')$} donc \deleted{définissent une catégorie torseur}
les classes d'isom d'objets font
un torseur sous \deleted{$\underline{\mathrm{Equ}}(C',T')/\inserted{\operatorname{Im}}\underline{\mathrm{Equ}}(C)$ par passage}
$\unsure{\mathrm{Bord}|\mathrm{Brd}}(C',T')$, donc définissent un torseur sous
$\mathrm{Bord}(C',T')/\operatorname{Im}\mathrm{Bord}(C)$. Le torseur ne dépend
plus du choix d'un $C$ dans la classe
de catégories envisagée\ldots\ Il faut donc expliciter
l'application composée

% B5 | diagram | foot, three aligned rows
\diagram{\begin{tikzcd}[ampersand replacement=\&,column sep=large,row sep=small]
\mathrm{Bord}(C)\arrow[r,"c"]\&\mathrm{Bord}(C,T)\arrow[r]\&\mathrm{Bord}(C',T')\\
\unsure{\beta|\rho}\arrow[r,mapsto]\&(\beta,u_\beta)\arrow[r,mapsto]\&(\beta',u_{\beta'})\\
{}_{O}P_{O}\arrow[u,phantom,"\unsure{\wr|\simeq}" description]\arrow[r,mapsto]\&N_{O/A}(P_O)\ni u_\beta\arrow[r,mapsfrom]\&N_{O'/A'}(P'_O)\ni(u_\beta)'
\end{tikzcd}}
% archive sheet number: 37

\scanpage{40}

% B1 | text | opening paragraph
Dans le cas qui nous intéresse ($\unsure{\xi|\zeta}\in\mathrm{Br}(\mathbb Q)$
d'indice $\unsure{n=2}$, ramifié \unsure{seulmt} en $2$ et $\infty$). \deleted{on trouve}
on trouve que $\unsure{\mathrm{Bord}|\mathrm{Brd}}(C)$ est la \deleted{\unsure{gr.}}\inserted{catégorie} de Picard
\unsure{splittée} d'invariants $\mathbb Z/2\mathbb Z$, $\mathbb Z/2\mathbb Z$, le générateur
$\unsure{\beta|\rho}$ de $\pi_0$ laquelle correspond au \deleted{\unsure{l'endomorphisme}} foncteur
qui à une courbe elliptique $E$ associe $E^{(p)}=E/\operatorname{Ker} F_E$)
correspond à l'unique idéal premier \inserted{$\mathfrak m$} de $O$ sur $p$,

% B2 | display | norm and choice in upper third
\[
\text{et on a}\qquad \underset{\substack{\Vert\\p\mathbb Z}}{N_{O/\mathbb Z}(\mathfrak m)}
\hookrightarrow N_{O/\mathbb Z}(O)=\mathbb Z,\qquad \text{on choisit }\bar\beta=p.
\]

% B3 | text | middle, first conditional
Ceci nous dit que \uline{si} $p$ \uline{est inversible dans} $A'$, alors
l'image de $(\beta,\bar\beta)$ dans $\mathrm{Brd}(C_{A'},T_{A'})\simeq\deleted{\unsure{\mathrm{Brd}(C_{A'})}}H^1(S',\mu_2)$
est la classe de $p$ modulo des carrés.

% B4 | text | middle, local case, with marginal notation
Par contre, prenons le cas où $A'=\mathbb Z_p$, alors
\deleted{on est dans un groupe}
$\mathrm{Brd}(C_p,T_p)\simeq\mathrm{Brd}(C_{A'},T_{A'})$ est une extension de $\mathbb Z/2\mathbb Z$ par
\margin{$C_p=C'$\newline $T_p=T'$}
$\mathbb Z_p^*/\mathbb Z_p^{*2}$ ($\simeq\mathbb Z/2\mathbb Z$ si $p\ne2$), en fait c'est
$\underbrace{\mathbb Z_p^*/\mathbb Z_p^{*2}}_{=H^1(\mathbb Z_p,\mu_2)\deleted{\unsure{\mathbb Z/(2)}}}\times
\underbrace{\mathbb Z/2\mathbb Z}_{\substack{\text{engendré par}\\(\mathfrak m,p)}}$, et les
classes d'isom \inserted{possibles des} données à l'$\infty$ \deleted{\unsure{liées aux}} forment un torseur
sous ce groupe (\deleted{\unsure{le plus souvent}} $\simeq\mathbb Z/2\mathbb Z\times\mathbb Z/2\mathbb Z$ si $p\ne2$).

% B5 | text | last paragraph, down to foot
D'autre part, si $A'=\mathbb R$, on peut considérer les classes d'isomorphie de la gerbe de lien $\mu_n$ sur $\mathbb R$ définie par $C$
(soit $\unsure{\mathscr G|\mathcal G}_{\mathbb R}$, où $\mathscr G$ est la gerbe sur $\mathbb Z[\frac1p]$) et \deleted{l'ensemble} des
classes d'\inserted{isomorphie d'}équivalences avec la gerbe-type sur $\mathbb R$
($=$ ensemble des \unsure{\uline{ple}ssions|\uline{pl}sections} de la gerbe), qui est un
torseur sous $H^1(\mathbb R,\mu_2)=\mathbb R^*/\mathbb R^{*2}$. On trouve donc
% archive sheet number: 39

\scanpage{42}

% B1 | text | opening paragraph and additions above lines
(Si on se donne la \deleted{\unsure{$\mathscr G$} \unsure{gerbe}} \inserted{catégorie de \unsure{Dessin}$_\otimes$ \unsure{polarisée}}
$(\unsure{C'|C'_f}(A'),T')$ sur $A'=\prod_{\ell\ne p,\infty}\mathbb Z_\ell$ comme
associée à la gerbe des racines carrées de $A'(1)$ (définie
par le choix d'une clôture algébrique $\bar{\mathbb F}_p$ de $\mathbb F_p$), \deleted{et} la
catégorie de Dessin$_\otimes$ $C(\mathbb Z_p)$ comme celle de $F$-cristaux sur $\bar{\mathbb F}_p$
(\unsure{par|pr} iso) qui sont de rang $2$, \deleted{\unsure{poids}} pente $\frac12$ \inserted{avec le $F$-cristal de Tate}, et la \deleted{\unsure{gerbe}}
catégorie $C(\mathbb Z_\infty)=C(\mathbb R)$ comme associée à la gerbe
type $\unsure{\mathscr G_\infty|\mathcal G_\infty}$ de groupe $\mu_2$ sur $\mathbb R$, on trouve en termes \deleted{\unsure{de}} $(C)$ un
certain torseur \inserted{$\mathcal T'(C)$} sous le groupe

% B2 | display | first product
\[
\deleted{\unsure{\mathbb Z_p}}\quad A'{}^*/A'{}^{*2}\times
(\mathbb Z_p^*/\mathbb Z_p^{*2}\times\mathbb Z/2\mathbb Z)
\times\mathbb R^*/\mathbb R^{*2}=\Gamma'
\]

% B3 | text | middle paragraph
Le groupe \inserted{$\mathbb Z/2\mathbb Z$} des classes d'isomorphie d'autoéquivalences de
\deleted{$C$} opère sur ce torseur, de façon compatible avec
son opération triviale sur lui-$=$, donc par translation,
et il opère donc via un hom \inserted{$f$ (\unsure{``Frobenius''})} de $\mathbb Z/2\mathbb Z$ dans ce groupe,
\deleted{\unsure{dont les}} l'image du \inserted{gén.} $\unsure{f}$ de $\mathbb Z/2\mathbb Z$ ds les 3 facteurs
étant respectivement

% B4 | display | middle lower, three components
\[
(p,\ \unsure{f|\beta}=(0,1),\ p)
\]

% B5 | text | small right-hand note linked to the last p
\margin{$\downarrow$ N.B. \deleted{\unsure{En fait là}, \unsure{pas} \unsure{légitime}},
\inserted{Comme $p\in\mathbb R^{*2}$} $\mathscr G_{\mathbb R}$ \unsure{orientable},
\unsure{entre} l'une des \unsure{ple}ssions,
\deleted{distinguer canoniqt \unsure{le}}
$\unsure{\mathrm{Pl}(C)}$ il y a un système transitif d'isomorphismes !}

% B6 | text | lower paragraph
Considérons le quotient

% B7 | display | lower centred quotient
\[
\Gamma=\Gamma'/f(\pm1)
\]

% B8 | text | last lines and addition above the last group
\deleted{On voit donc}
et soit $\mathcal T(C)$ le quotient de $\mathcal T'(C)$ par $\pm1$, qui
est donc un torseur sous $\Gamma$. Ce torseur, à isom unique
près, ne dépend pas de $C$, \inserted{on peut l'appeler $\mathcal T$ ou $\mathcal T(\bar{\mathbb F}_p)$.} D'ailleurs si $\Gamma'_0$ est le
ss-groupe de $\Gamma'$\quad $A'{}^*/A'{}^{*2}\times\mathbb Z_p^*/\mathbb Z_p^{*2}\times\mathbb R^*/\mathbb R^{*2}
=H^1_{\deleted{\unsure{\mathrm{ad}}}}(\mathbb A,\mu_2)$ \inserted{adèles}
on trouve que $\Gamma'\simeq\Gamma'_0\times f(\pm1)$, donc $\Gamma\simeq\Gamma'_0$ can, donc
% archive sheet number: 41

\scanpage{44}

% B1 | text | top line and equality beneath right end
$\mathcal T(\bar{\mathbb F}_p)$ est un torseur canonique sous $H^1(\mathbb A,\mu_2)$

% B2 | display | top right
\[
\underset{\substack{\Vert\\\prod_\ell\mathbb Z_\ell^*/\mathbb Z_\ell^{*2}\\
\text{(y inclus $p$ et $\infty$)}}}{\vphantom{H^1(\mathbb A,\mu_2)}}
\]

% B3 | text | upper left, vertical note with brace
% margin note written sideways
\margin{\deleted{\unsure{à isom près}} On précisera l'objet : \unsure{quintuplet|quaterne} $(\unsure{\mathscr E}M,o)$ \uline{à isom près} (il y a deux classes d'isomorphie possibles si $(M,o)$ est l'une d'eux, l'autre est $(M^{(p)},o')$ que l'on en déduit en appliquant frobenius}

% B4 | text | main upper paragraph
Ce torseur à priori dépend encore de $\bar{\mathbb F}_p$, \deleted{\unsure{mais il faut}} et cette
dépendance est précisée par la connaissance de l'opération
de $\hat{\mathbb Z}\ni\unsure{f|F}$ sur $\mathcal T(\bar{\mathbb F}_p)$. Cette opération se fait via
les opérations de $\hat{\mathbb Z}$ sur $C(A',T')$, \deleted{\unsure{$E$ étant tel}} \inserted{(\unsure{cycles|opér.} habituelles sur $A'(1)$, $f$ opère par $p$) et} $C(A_p,T_p)$
(via opération de $\hat{\mathbb Z}f$ sur $C(A_p,T_p)$ par frobenius) et
l'opération triviale sur $C(A_\infty,T_\infty)$. Bon on trouve que
$f\in\hat{\mathbb Z}$ opère par le $=$ élément $(p,f,p)$ que plus
haut, donc

% B5 | display | central isolated line
\[
f\text{ opère trivialement sur }\mathcal T(\bar{\mathbb F}_p)
\]

% B6 | text | central conclusion
Ainsi, on a un torseur canonique \inserted{$\mathcal T$} sous $H^1(\mathbb A,\mu_2)$.

% B7 | text | boxed mark and NB in lower half
$\square$ \uline{NB} Ceci suggère qu'on a une \deleted{section canonique} \inserted{gerbe triviale (au}
\unsure{préhistorie|histoire}) de \deleted{\unsure{faisceaux}} $\mu_2$ sur $\mathbb A$, dont $\mathcal T$ serait
le torseur des classes d'isomorphie. On peut
en effet considérer les couples $(C,\deleted{\unsure{f}}\alpha)$, où
$C$ est une catégorie de Dessin$_\otimes$ de type $\unsure{\xi|\zeta}\in\mathrm{Br}(\mathbb Q)$, \inserted{d'où $T_C=T$,} et
$\alpha$ une \deleted{\unsure{équiv}} classe d'équivalence de classes d'isomorphie
de \inserted{l'équival.} $(C,T)_{\mathbb Z_p}$ avec $(C(\mathbb Z_p),T(\mathbb Z_p))$, modulo équivalences
provenant de $\underline{\underline{\mathrm{Pic}}}(\mathbb Z_p)$ ; i.e. une section du torseur de
groupe $\mathbb Z/2\mathbb Z$ déduit de celui de groupe $\underline{\mathrm{Equ}}(C(\mathbb Z_p),T(\mathbb Z_p))=\mathbb Z_p^*/\mathbb Z_p^{*2}\times\mathbb Z/2\mathbb Z$
déduit par extension du gpe structural à l'aide de $\mathrm{pr}_2$.

% B8 | text | lower left, sloping note connected to side note
% slanted
\margin{\unsure{trop} compliqué il suffit des \unsure{fichues|fichus} classes d'isom d'\unsure{isomorph.} d'équivalences de cat. de Dessin$_\otimes$ :
$C_{\mathbb Z_p}\simeq C(\mathbb Z_p)$}

% B9 | text | lower left, second note and cancellation
% slanted
\margin{en termes concrets : un $\alpha$ est donné (si $C$ défini par $O$) en \unsure{se} donnant une struct. de $O\otimes_{\mathbb Z}\mathbb Z_p$ sur l'anneau des endos $M_0$ ; d'un cristal sur $\bar{\mathbb F}_p$ (de rang $2$,
\deleted{pente $\frac12$ \unsure{et} \unsure{$\bigwedge^2$} \unsure{$A$} / \unsure{un choix d'un lien de Tate}}}

% B10 | text | bottom edge
Alors tout $C$ admet un $\alpha$, et deux $(C,\alpha),(C',\alpha')$
% archive sheet number: 43

\scanpage{46}

% B1 | text | cancelled opening, with addition above it
\deleted{sont \uline{\unsure{canoniquement}} équivalentes, \inserted{l'équivalence devenant \unsure{à} isom près} Si on \unsure{veut} les deux
essentiellement identifier entre eux les différents $(C,\alpha)$,
on trouve ainsi une catégorie de \uline{Dessin}$_\otimes$
\uline{canonique}}

% B2 | text | upper paragraph, continuation after cancellation
sont équivalents : la catégorie des équivalences
de $(C,\alpha)$ avec $(C',\alpha')$ est un groupoïde connexe de
$\pi_1$ $\mathbb Z/2\mathbb Z$ : i.e. \unsure{$\alpha$ donne|toutes} les équivalences sont
isom, et si on a deux équiv. $\varphi,\varphi'$, il y a \uline{à un
signe près} un seul isom. entre $\varphi$ et $\varphi'$.

% B3 | text | main middle paragraph
Si on se donne maintenant un couple $(C,\alpha)$,
il lui est associé une \deleted{torseur sous} gerbe sur $\mathbb A=\prod_{\substack{\ell\ (\text{y compris}\\\ell=p\text{ et }\ell=\infty)}}\mathbb Z_\ell$
\inserted{de lien $\mu_2$,} savoir celle des équivalences \inserted{de $(C,T)$} avec le \deleted{\unsure{gerbe}} couple
type de $(C(\mathbb A),T)$, qui en la place $p$ soit
compatible avec la donnée de $\alpha$. Cette gerbe
$\unsure{\mathscr G|\mathcal G}_{(C,\alpha)}$ est déterminée, ``à un signe près'' \deleted{\unsure{ou}} équivalence
unique près, indépendamment du choix de $(C,\alpha)$
(puisqu'il en est ainsi de $(C,\alpha)$ lui-$=$\ldots), \deleted{et}
$\deleted{\hat{\mathbb Z}=\mathrm{Gal}(\bar{\mathbb F}_p/\mathbb F_p)}$ \deleted{y opère trivialmt.} mais elle dépend
a priori, \inserted{$=$ pour $(C,\alpha)$ fixé,} du choix d'un $\bar{\mathbb F}_p$ \deleted{(ou du moins
de $\mathbb F_{p^2}$)} -- tout comme la structure $\unsure{-|\sim}$ elle-$=$ !

% B4 | text | lower left, struck marginal note
% slanted
\margin{\deleted{car le $F$-cristal de rang $2$, pente $\frac12$ sur $\bar{\mathbb F}_p$ se \unsure{descendait|descendant} \uline{canoniquement} à $\mathbb F_{p^2}$}}

% B5 | text | last paragraph
On peut s'en tirer en définissant \uline{canoniqt} un
$\mathbb F_{p^2}$ (ou $=$ un $\bar{\mathbb F}_p$, si on y tenait\ldots) par les polynomes
cyclotomiques $\Phi_{p^n-1}$,

% B6 | display | foot
\[
\bar{\mathbb F}_p=\varinjlim_n\left(\mathbb F_{p^n}=\mathbb F_p[T]/\Phi_{p^n-1}(T)\right)
\qquad \ldots\ldots
\]
\unsure{;}
% archive sheet number: 45

\scanpage{48}

% B1 | text | upper paragraph
Considérons cette gerbe $\unsure{\mathscr G|\mathcal G}$ \deleted{\unsure{sur}} sur $\mathbb A=\prod_p\mathbb Z_p$, de
groupe $\mu_2$, définie ``au signe près'' (de
façon précise, on a une $2$-catégorie $\unsure{\mathscr D|\mathcal D}$ des $(C,\alpha)$,
qui satisfait $\pi_0=\pi_1=0$, $\pi_2=\mathbb Z/2\mathbb Z$, et \deleted{que} on
a \deleted{\unsure{une}} \inserted{une famille de Gerbes} indexée par $\mathscr D$, i.e. un $2$ foncteur
de $\mathscr D$ dans la $2$-catégorie des \deleted{$2$} gerbes sur $\mathbb A$,
groupe $\mu_2$). Dans cette gerbe, on a une
``\uline{section définie à isom can. au signe près}''. $]]$

% B2 | text | faint left margin opposite underlining
% slanted
\margin{attention : \unsure{réfléchir|révérifier}, s'il n'y a pas un deuxième signe !}

% B3 | text | central separator and new paragraph
\noindent\rule{2cm}{0.4pt}
Autre façon d'\unsure{épingler} : le \deleted{catégorie des équ} \inserted{foncteur des}
catégories de Picard \unsure{composé}

% B4 | diagram | middle, two horizontal arrows and isomorphism beneath the target
\diagram{\begin{tikzcd}[ampersand replacement=\&,column sep=small,row sep=tiny]
\underline{\mathrm{Equ}}(C_{\unsure{\otimes|n}})\arrow[r]\&\underline{\mathrm{Equ}}(C,T)\arrow[r]\&\underline{\mathrm{Equ}}(C_p,T_p)\arrow[d,phantom,"{\wr\!\mid\;\mathrm{can}}" description]\\
\&\&\mathrm{Equ}(C(\mathbb Z_p),T(\mathbb Z_p))
\end{tikzcd}}

% B5 | text | lower half through crossed lines
est pleinement fidèle, la catégorie de Picard quotient est
donc une catégorie rigide, qui n'est autre que
$[0\mapsto\mathbb Z_p^*/\mathbb Z_p^{*2}]$. Donc la catégorie des équivalences de
$(C,T)_{\mathbb Z_p}$ avec $(C(\mathbb Z_p),T(\mathbb Z_p))$ \deleted{\unsure{$\to$}} \uline{modulo} \uline{action de} $\underline{\mathrm{Equ}}(C)$
est une catégorie rigide, \deleted{torseur} dont l'ens des
classes d'isom. est un $\mathcal T_p$, torseur canonique sous
$\mathbb Z_p^*/\mathbb Z_p^{*2}$ (on choisira \deleted{lequel}\inserted{qui est}!!!), indépendant à isom
can. près de $C$. \deleted{On a un \unsure{choix} can} \inserted{pour \deleted{\unsure{quel}} $\xi_p\in\mathcal T_p$.}
\deleted{Grâce à cet élément, \unsure{le choix} d'une \unsure{équiv}}
\inserted{on peut considérer} les
couples $(C,\alpha)$, où \deleted{\unsure{$C$}} $\alpha$ est une équivalence
$(C,T)_{\mathbb Z_p}\simeq(C(\mathbb Z_p),T(\mathbb Z_p))$ \uline{dans la classe} \uline{de} $\xi_p$, \deleted{et}
\deleted{\unsure{on}} Si on a deux tels couples, ils sont
\deleted{isomorphes} équiv. \uline{de} \uuline{façon canonique}
% archive sheet number: 47

\scanpage{50}

% B1 | text | two opening lines
Donc on peut \deleted{considérer que} \inserted{identifier} les $(C,\alpha)$ \deleted{\unsure{sont}} entre
eux \deleted{\unsure{à deux}}, on trouve donc

% B2 | display | upper centre
\[
[C(\xi),\alpha(\xi)]\qquad\text{canonique.}
\]

% B3 | text | boxed and crossed opening of abandoned paragraph
\deleted{On peut \unsure{donc|alors} considérer sur $\mathbb A''$\inserted{$=\mathbb A'\times A_\infty$} (anneau des adèles
premières à $p$), \deleted{\unsure{un}} la gerbe $g''$ \inserted{de groupe $\mu_2$} qui sur $A'$ est la
gerbe des racines carrées de $A'(1)$, et sur $A_\infty$ est la
gerbe type $g'_\infty$ déjà envisagée, alors}

% B4 | text | main paragraph, upper middle
En fonction de ce choix, on trouve une gerbe
canonique $\unsure{\mathscr G|\mathcal G}(\xi)\deleted{\unsure{=}}$ sur l'anneau des adèles,
de groupe $\mu_2$, savoir la gerbe qui sur $A''=A'\times A_\infty$ est
la gerbe des équivalences \deleted{\unsure{en}} de la gerbe \deleted{$\mathscr G\unsure{(\mathbb Z_p)}$} \inserted{de gpe $\mu_2$}
définie sur $\mathbb Z[\frac1p]$ par $(C,T)$ \deleted{\unsure{avec la}} avec la gerbe
$g'_{A''}$ déjà envisagée plus haut (sur $A'$ c'est celles
des \unsure{$\sqrt{\ }$} de $A'(1)$, sur $A_\infty$ c'est $g'_\infty$ déjà envisagée
plus haut) et qui sur $A_p$ est \uuline{la} \uline{gerbe} \deleted{des}\inserted{\uline{triviale}}.

% B5 | text | small middle-lower writing, then a boxed deletion
Cette gerbe est \deleted{\unsure{la}} triviale, et l'ens des \deleted{classes}\inserted{sections} à isom
\deleted{\unsure{équiv. des}} près est \deleted{le} torseur $\mathcal T$ sous $H^1(\mathbb A,\mu_2)$ déjà
\inserted{envisagé plus haut, qui est \unsure{ess.} indép. de $\xi$.}
\deleted{définie par les équiv \unsure{$(-)$} dit, \unsure{on aurait} une
\uline{section canonique} \deleted{\unsure{$\mathscr G$}} \unsure{$s(\xi)$} de $\mathscr G(\xi)$}

% B6 | text | far left opposite large cancellation
\margin{\deleted{\unsure{D'autre part la}}}

% B7 | text | bottom paragraph
Tout ceci posé, la théorie \inserted{de Deuring} des courbes elliptiques
d'inv. de Hasse nul nous fournit

% B8 | text | bottom item a
a) un élément $\xi_p\in\mathcal T_p$ (l'\uline{élément de Deuring})

% B9 | text | bottom item b
b) une section \deleted{correspondante} $\bar\xi$ de $\mathscr G(\xi_p)$, définie
à isom près, unique \uline{au signe près à l'infini},
dont la valeur en $p$ \deleted{\unsure{a comme}} est la section unité
\deleted{\unsure{can.} \unsure{unique}} de $\mathscr G(\xi_p)_{\mathbb Z_p}$, et la valeur en $\infty$ est la
polarisation donnée par $T$\ldots

% B10 | text | foot, concluding lines
On trouve donc une \deleted{section}\inserted{élémt} $\{\xi\}$ de $\mathcal T$, dont la
comp. suivant $p$ est $\xi_p$.
% archive sheet number: 49

\scanpage{52}

% B1 | text | opening lines
Á l'infini, la situation n'est pas très belle encore. Il
vaut mieux dire :

% B2 | text | first full paragraph
Pour tt $\xi_p\in\mathcal T_p$, \deleted{la valeur en $\infty$ de $\mathscr G(\xi)$ est} \inserted{considérer la gerbe $\unsure{h|\mathfrak h}(\xi_p)$} de
groupe $\mu_2$ sur $\unsure{\mathbb Z_{1/p}\mathbb Z|\mathbb Z[1/p]}$ définie par $(C(\xi_p),T(\xi_p))$, \deleted{En fait} \inserted{Considérons}
$\deleted{\unsure{\mathscr G}} h(\xi_p)_{\mathbb R}$, qui est une gerbe \uline{canonique} sur $\mathbb R$ de groupe
$\mu_2$. \deleted{J'ignore si} \inserted{pour $\xi_p=$ élément de Deuring,} elle a une description indépendante
de la théorie des catégories de Deuring (et de celle
des courbes elliptiques). \unsure{Ce qu'on|Ce que} \unsure{sait|sont} en tous cas
qu'elle est non triviale ; elle est, pour $\xi_p=$ élément de Deuring,
munie d'une polarisation $\xi_\infty$ (ce qui revient à la
donnée d'une classe d'isomorphie \inserted{d'équivalences} avec la gerbe
type $g_\infty$, \uline{mais il reste} \uline{l'ambiguïté du signe}.

% B3 | text | middle, category over the reals
Catégorie $C(\xi_p)_\infty$ \deleted{\unsure{est}} $\simeq$ catégorie des vectoriels \inserted{à droite $V$ de dim $1$} \deleted{\unsure{qui}} sur
$\mathbb H$ ($=$ quaternions ordinaires), avec
\deleted{\unsure{leur} $V$} foncteur norme
réduite vers cat. de Picard de $\mathbb R$.

% B4 | text | lower middle, functor
Ainsi il y a un foncteur ``canonique au signe près''

% B5 | diagram | lower middle, arrow and lower additions
\diagram{\begin{tikzcd}[ampersand replacement=\&,column sep=large]
C(\xi_p)\arrow[r,"{\deleted{\unsure{\xi}}c_\infty}"]\&\deleted{\unsure{C(\mathbb R)}\ \text{catégorie}}\,C(\mathbb R)\\
(T(\xi_p)\arrow[r,draw=none]\&\deleted{\unsure{\xi}}\;\mathbb R)
\end{tikzcd}}

% B6 | text | left-hand leader and right-hand gloss
\margin{$\longleftarrow$ n'est pas spécial à $C(\xi_p)$, vaut \deleted{\unsure{pour}} pour tt $C$}
\margin{la \unsure{factorisation|fonction} \unsure{Arr.} par les vectls quaternioniens de dim $1$ \uuline{unitaires}}

% B7 | text | struck lines below diagram, with continuation
\deleted{et qui associe à chaque élément de
\unsure{la} \unsure{leur} \unsure{définit} un autre}
tout autre choix étant isom
à ça, l'isom \uline{défini au signe près} (\unsure{chic|chicos} !).

% B8 | text | foot, last paragraph
Prenant $\xi_p=$ élément de \deleted{\unsure{Eisenstein}}\inserted{Deuring}, on obtient un
foncteur \inserted{ess. canonique} des courbes elliptiques de Hasse nul
vers les vectoriels \inserted{\unsure{quatrs} unitaires} de dim $1$. La \deleted{\unsure{transp}} Norme
correspond à la polarisation, la \unsure{transposition|transpiration}, la transposition
% archive sheet number: 51
